\section{Conserved Currents and Charges}
\label{sec:conserved_currents_charges}

Noether's theorem produces a local equation,

\[
\partial_{\mu}j^{\mu}=0.
\]

To interpret this equation physically, it must be separated into a density and a spatial flux. The integrated density is conserved only when the flux through the boundary is controlled.

\subsection{Continuity equation}

Use

\[
x^{0}=ct,
\qquad
\partial_{0}
=
\frac{1}{c}
\frac{\partial}{\partial t}.
\]

Write the four-current as

\[
\boxed{
j^{\mu}
=
\left(
c\rho,
\mathbf{j}
\right)
}.
\]

Then

\[
\partial_{\mu}j^{\mu}
=
\frac{1}{c}
\frac{\partial}{\partial t}
\left(
c\rho
\right)
+
\nabla\cdot\mathbf{j}.
\]

Therefore

\[
\boxed{
\frac{\partial\rho}{\partial t}
+
\nabla\cdot\mathbf{j}
=
0
}.
\]

The quantity \(\rho\) is the local charge density, and \(\mathbf{j}\) is the spatial current or flux density.

\subsection{Integrated charge}

For a spatial region \(V\), define

\[
\boxed{
Q_V(t)
=
\int_V
\rho(\mathbf{x},t)
\,\mathrm{d}^{3}x
}.
\]

Differentiate:

\[
\frac{\mathrm{d}Q_V}{\mathrm{d}t}
=
\int_V
\frac{\partial\rho}{\partial t}
\,\mathrm{d}^{3}x.
\]

Using the continuity equation,

\[
\frac{\mathrm{d}Q_V}{\mathrm{d}t}
=
-
\int_V
\nabla\cdot\mathbf{j}
\,\mathrm{d}^{3}x.
\]

The divergence theorem gives

\[
\boxed{
\frac{\mathrm{d}Q_V}{\mathrm{d}t}
=
-
\oint_{\partial V}
\mathbf{j}\cdot\mathbf{n}
\,\mathrm{d}A
}.
\]

The charge inside \(V\) changes only because current crosses the boundary.

\subsection{Conservation on all space}

For

\[
V=\mathbb{R}^{3},
\]

the total charge is

\[
\boxed{
Q
=
\int_{\mathbb{R}^{3}}
\rho
\,\mathrm{d}^{3}x
=
\frac{1}{c}
\int_{\mathbb{R}^{3}}
j^{0}
\,\mathrm{d}^{3}x
}.
\]

If the fields decay sufficiently rapidly that

\[
\lim_{R\to\infty}
\oint_{|\mathbf{x}|=R}
\mathbf{j}\cdot\mathbf{n}
\,\mathrm{d}A
=
0,
\]

then

\[
\boxed{
\frac{\mathrm{d}Q}{\mathrm{d}t}=0
}.
\]

Local conservation alone is not enough. Boundary behavior is part of the argument.

\subsection{Finite regions and balance laws}

In a laboratory region, the charge need not be constant. The correct statement is the balance law

\[
Q_V(t_2)-Q_V(t_1)
=
-
\int_{t_1}^{t_2}
\oint_{\partial V}
\mathbf{j}\cdot\mathbf{n}
\,\mathrm{d}A\,\mathrm{d}t.
\]

A decrease of charge inside the region is exactly balanced by outward flux.

This interpretation applies not only to electric charge but also to energy, momentum, and other Noether quantities.

\subsection{Boundary conditions}

Several common conditions make the boundary flux vanish:

\begin{enumerate}
    \item compact support of the current;
    \item sufficiently rapid decay at spatial infinity;
    \item periodic boundary conditions, for which opposite-face fluxes cancel;
    \item reflecting boundaries with
    \[
    \mathbf{j}\cdot\mathbf{n}=0;
    \]
    \item integration over a closed system whose boundary is not crossed by the relevant current.
\end{enumerate}

The appropriate condition depends on the physical problem.

\subsection{Current nonuniqueness and charge}

Let

\[
j'^{\mu}
=
j^{\mu}
+
\partial_{\nu}B^{\mu\nu},
\qquad
B^{\mu\nu}
=
-B^{\nu\mu}.
\]

The charge density changes by

\[
\rho'
=
\rho
+
\frac{1}{c}
\partial_iB^{0i}.
\]

Therefore

\[
Q'-Q
=
\frac{1}{c}
\int_{\mathbb{R}^{3}}
\partial_iB^{0i}
\,\mathrm{d}^{3}x
=
\frac{1}{c}
\oint_{S_{\infty}}
B^{0i}n_i
\,\mathrm{d}A.
\]

If the improvement term decays sufficiently rapidly, then

\[
Q'=Q.
\]

Thus different local representatives of a current can define the same physical charge.

\subsection{Worked example: shift charge}

For the massless scalar shift symmetry,

\[
j^{\mu}
=
\partial^{\mu}\phi.
\]

Since

\[
\partial^{0}\phi
=
\frac{1}{c}
\frac{\partial\phi}{\partial t},
\]

the density is

\[
\rho
=
\frac{j^{0}}{c}
=
\frac{1}{c^{2}}
\frac{\partial\phi}{\partial t}.
\]

The charge is

\[
\boxed{
Q_{\mathrm{shift}}
=
\frac{1}{c^{2}}
\int
\frac{\partial\phi}{\partial t}
\,\mathrm{d}^{3}x
}.
\]

Its time derivative is

\[
\frac{\mathrm{d}Q_{\mathrm{shift}}}{\mathrm{d}t}
=
\frac{1}{c^{2}}
\int
\frac{\partial^{2}\phi}{\partial t^{2}}
\,\mathrm{d}^{3}x.
\]

Using the massless wave equation,

\[
\frac{1}{c^{2}}\phi_{tt}
=
\nabla^{2}\phi,
\]

we find

\[
\frac{\mathrm{d}Q_{\mathrm{shift}}}{\mathrm{d}t}
=
\int
\nabla^{2}\phi
\,\mathrm{d}^{3}x
=
\oint
\nabla\phi\cdot\mathbf{n}
\,\mathrm{d}A.
\]

The charge is constant when the boundary term vanishes.

\subsection{Local conservation versus static density}

A conserved current does not imply

\[
\frac{\partial\rho}{\partial t}=0
\]

at every point. The density can move through space. The continuity equation states that local changes are balanced by flux:

\[
\frac{\partial\rho}{\partial t}
=
-\nabla\cdot\mathbf{j}.
\]

A wave packet can carry a conserved total charge while its density profile changes position and shape.

\subsection{Charge as generator}

In Hamiltonian formulations, a conserved Noether charge often generates the corresponding infinitesimal transformation through Poisson brackets. Schematically,

\[
\delta\phi_a
=
\varepsilon
\left\{
\phi_a,Q
\right\}.
\]

This deeper relation explains why conserved charges encode symmetry operations. A complete Hamiltonian treatment lies beyond the present derivation, but the structural connection is worth recording.

\subsection{Diagnostic checks}

For a proposed current:

\begin{enumerate}
    \item compute \(\partial_{\mu}j^{\mu}\) directly;
    \item use the field equations only after the divergence is expanded;
    \item identify the density \(j^{0}/c\);
    \item identify the spatial flux \(\mathbf{j}\);
    \item check the dimensions of the integrated charge;
    \item state the boundary condition needed for \(\mathrm{d}Q/\mathrm{d}t=0\).
\end{enumerate}

\subsection{Common mistakes}

Typical errors include:

\begin{enumerate}
    \item treating \(\partial_{\mu}j^{\mu}=0\) as if it implied a time-independent density;
    \item defining a charge without the factor appropriate to the convention \(x^{0}=ct\);
    \item discarding the boundary flux without stating a decay or boundary condition;
    \item confusing the current \(j^{\mu}\) with its integrated charge \(Q\);
    \item concluding that two improved currents define different charges without checking the surface term;
    \item using a global conservation statement when only a finite-region balance law is justified.
\end{enumerate}

\subsection{What has been established}

A conserved four-current produces the continuity equation

\[
\frac{\partial\rho}{\partial t}
+
\nabla\cdot\mathbf{j}
=
0.
\]

The associated charge is constant only when the boundary flux vanishes. This local-to-global passage will now be applied to spacetime translations, producing energy and momentum conservation.
